\documentclass[addpoints]{exam}
\usepackage{cmbright}
\usepackage{amsmath}
\usepackage{amssymb}
\usepackage{graphicx}
\usepackage{url}
\usepackage{mhchem}
\setlength\fillinlinelength{2em}

\header{CHE 525: math for CBEE's}{homework \#5}{fall 2025}

\begin{document}

do all of these questions by hand, i.e. without a computer. of course, check your answers with Julia.

\begin{questions}
\subsection*{orthonormal bases}

%4.4 p5
\question find two orthogonal vectors in the plane $x+y+2z=0$. make them orthonormal.

\question if the matrix $Q$ has orthonormal columns, what is the least squares solution $\hat{\mathbf{x}}$ to $Q\mathbf{x}=\mathbf{b}$? 
%what is the projection of $\mathbf{b}$ onto the column space of $Q$?

% 4.4 P15
\question 
	\begin{parts}
\part find orthonormal vectors $\mathbf{q}_1$, $\mathbf{q}_2$ that span the column space of:
\begin{equation}
A =
\begin{bmatrix}
1 & 1\\
2 & -1\\
-2 & 4
\end{bmatrix}
\end{equation}

\part write the $Q=[\mathbf{q}_1, \mathbf{q}_2]$ and $R$ matrices for the factorization $A=QR$. (check that it works in Julia by \texttt{Q * R}. note \texttt{Q, R = qr(A)} should give the same directions but maybe different signs.)

\part faced with the overdetermined system $A\mathbf{x}=\begin{bmatrix}2 \\ 2 \\ 7\end{bmatrix}$, use your $A=QR$ factorization to find the least squares solution $\hat{\mathbf{x}}$. you should arrive at a 2 by 2 triangular system to solve for $\hat{\mathbf{x}}$. (check your answer with \texttt{A \textbackslash b} in Julia.) 

\part finally, find a third vector $\mathbf{q}_3$ that makes $\mathbf{q}_1, \mathbf{q}_2, \mathbf{q}_{3}$ an orthonormal basis for $\mathbb{R}^{3}$. which of the four fundamental subspaces of $A$ contains $\mathbf{q}_{3}$?


\end{parts}
\end{questions}
 (\emph{source}: ``Introduction to Linear Algebra'' by Strang, 5th Ed., 2016)
\end{document}
